% Options for packages loaded elsewhere
\PassOptionsToPackage{unicode}{hyperref}
\PassOptionsToPackage{hyphens}{url}
\documentclass[
  ignorenonframetext,
]{beamer}
\newif\ifbibliography
\usepackage{pgfpages}
\setbeamertemplate{caption}[numbered]
\setbeamertemplate{caption label separator}{: }
\setbeamercolor{caption name}{fg=normal text.fg}
\beamertemplatenavigationsymbolsempty
% remove section numbering
\setbeamertemplate{part page}{
  \centering
  \begin{beamercolorbox}[sep=16pt,center]{part title}
    \usebeamerfont{part title}\insertpart\par
  \end{beamercolorbox}
}
\setbeamertemplate{section page}{
  \centering
  \begin{beamercolorbox}[sep=12pt,center]{section title}
    \usebeamerfont{section title}\insertsection\par
  \end{beamercolorbox}
}
\setbeamertemplate{subsection page}{
  \centering
  \begin{beamercolorbox}[sep=8pt,center]{subsection title}
    \usebeamerfont{subsection title}\insertsubsection\par
  \end{beamercolorbox}
}
% Prevent slide breaks in the middle of a paragraph
\widowpenalties 1 10000
\raggedbottom
\AtBeginPart{
  \frame{\partpage}
}
\AtBeginSection{
  \ifbibliography
  \else
    \frame{\sectionpage}
  \fi
}
\AtBeginSubsection{
  \frame{\subsectionpage}
}
\usepackage{iftex}
\ifPDFTeX
  \usepackage[T1]{fontenc}
  \usepackage[utf8]{inputenc}
  \usepackage{textcomp} % provide euro and other symbols
\else % if luatex or xetex
  \usepackage{unicode-math} % this also loads fontspec
  \defaultfontfeatures{Scale=MatchLowercase}
  \defaultfontfeatures[\rmfamily]{Ligatures=TeX,Scale=1}
\fi
\usepackage{lmodern}
\ifPDFTeX\else
  % xetex/luatex font selection
\fi
% Use upquote if available, for straight quotes in verbatim environments
\IfFileExists{upquote.sty}{\usepackage{upquote}}{}
\IfFileExists{microtype.sty}{% use microtype if available
  \usepackage[]{microtype}
  \UseMicrotypeSet[protrusion]{basicmath} % disable protrusion for tt fonts
}{}
\makeatletter
\@ifundefined{KOMAClassName}{% if non-KOMA class
  \IfFileExists{parskip.sty}{%
    \usepackage{parskip}
  }{% else
    \setlength{\parindent}{0pt}
    \setlength{\parskip}{6pt plus 2pt minus 1pt}}
}{% if KOMA class
  \KOMAoptions{parskip=half}}
\makeatother
\usepackage{longtable,booktabs,array}
\newcounter{none} % for unnumbered tables
\usepackage{calc} % for calculating minipage widths
\usepackage{caption}
% Make caption package work with longtable
\makeatletter
\def\fnum@table{\tablename~\thetable}
\makeatother
\setlength{\emergencystretch}{3em} % prevent overfull lines
\providecommand{\tightlist}{%
  \setlength{\itemsep}{0pt}\setlength{\parskip}{0pt}}
% Slide layout defaults for warning-clean scientific decks.
\usepackage{etoolbox}
\IfFileExists{xurl.sty}{\usepackage{xurl}}{}
\IfFileExists{seqsplit.sty}{\usepackage{seqsplit}}{\newcommand{\seqsplit}[1]{#1}}
\protected\def\breakseq#1{\seqsplit{#1}}
\protected\def\breaktt#1{\begingroup\ttfamily\seqsplit{#1}\endgroup}
\setlength{\emergencystretch}{6em}
\tolerance=5000
\hbadness=10000
\hfuzz=1pt
\setlength{\tabcolsep}{2pt}
\AtBeginEnvironment{longtable}{\tiny\renewcommand{\arraystretch}{0.86}\setlength{\tabcolsep}{1pt}}
\AtBeginEnvironment{tabular}{\tiny\renewcommand{\arraystretch}{0.86}\setlength{\tabcolsep}{1pt}}

% Natbib and cross-reference fallbacks — slides don't load natbib
% or cleveref, but combined-PDF manuscript prose may emit these
% commands. The fallback renders citations as a bracketed key list
% and cross-references as detokenized labels so slides don't fail on
% undefined control sequences or raw underscores. \providecommand is
% a no-op if packages load later.
\providecommand{\citep}[1]{[#1]}
\providecommand{\citet}[1]{#1}
\providecommand{\citealp}[1]{#1}
\providecommand{\citeauthor}[1]{#1}
\providecommand{\citeyear}[1]{#1}
\providecommand{\cref}[1]{\texttt{\detokenize{#1}}}
\providecommand{\Cref}[1]{\texttt{\detokenize{#1}}}

\newtheorem{proposition}{Proposition}
\newtheorem{hypothesis}{Hypothesis}
\usepackage{bookmark}
\IfFileExists{xurl.sty}{\usepackage{xurl}}{} % add URL line breaks if available
\urlstyle{same}
% fallback for those not using the hyperref driver hyperxmp:
\makeatletter
\@ifundefined{xmpquote}{\newcommand{\xmpquote}[1]{#1}}{}
\makeatother
\hypersetup{
  hidelinks,
  pdfcreator={LaTeX via pandoc}}

\author{\texorpdfstring{}{}}
\date{}

\begin{document}

\section{Results}\label{results}

\begin{frame}[fragile,allowframebreaks]{Fixture Release Packet}
\protect\phantomsection\label{fixture-release-packet}
The fixture release packet contains fourteen segments spanning four
classification levels: UNCLASSIFIED (ten segments), CUI (one segment),
SECRET (two segments), and TOP\_SECRET (one segment). Four segments
carry source controls (drawn from the HUMINT, SIGINT, and IMINT
disciplines). Twenty-one redaction decisions are applied across those
four segments, using all five bounded reasons:
\texttt{source\_identity}, \breaktt{operational\_detail},
\breaktt{time\_place\_selector}, \texttt{legal\_privilege}, and
\texttt{privacy}.

The audit produces:

\begin{itemize}
\tightlist
\item
  \textbf{Releasable}: true (no error-level findings after redaction).
\item
  \textbf{Release safety score}: weighted by error count, warning count,
  and mosaic risk.
\item
  \textbf{Redaction coverage}: 1.0 (all sensitive segments have at least
  one decision).
\item
  \textbf{Mosaic risk score}: residual markers normalized by segment and
  pattern count.
\item
  \textbf{Findings}: warning-level findings for residual markers in
  sanitized text.
\end{itemize}
\end{frame}

\begin{frame}[fragile,allowframebreaks]{Source-Safe Redaction Ledger}
\protect\phantomsection\label{source-safe-redaction-ledger}
The redaction ledger records each decision with:

\begin{itemize}
\tightlist
\item
  A decision ID derived from the SHA-256 of the segment ID, span,
  reason, and replacement.
\item
  The segment ID, start, end, span length, reason, and replacement.
\item
  A \texttt{valid\_span} flag indicating whether the span falls within
  the segment text.
\item
  A \breaktt{source\_span\_sha256} hash of the redacted text---present
  only for valid spans, absent for invalid or orphan decisions.
\end{itemize}

The ledger never exposes source text. It provides reproducible audit
evidence without disclosure risk.
\end{frame}

\begin{frame}[fragile,allowframebreaks]{Segment Hash Manifest}
\protect\phantomsection\label{segment-hash-manifest}
The hash manifest records, for each segment:

\begin{itemize}
\tightlist
\item
  \texttt{source\_sha256}: SHA-256 of the original segment text.
\item
  \texttt{public\_sha256}: SHA-256 of the redacted segment text.
\end{itemize}

This allows downstream verification that the public release matches the
audited version without comparing source text directly.
\end{frame}

\begin{frame}[fragile,allowframebreaks]{Review Gate}
\protect\phantomsection\label{review-gate}
The release gate requires three reviewer roles: originator,
classification reviewer, and release authority. Each reviewer provides a
non-empty rationale. The fixture reviews all approve, yielding:

\begin{itemize}
\tightlist
\item
  \textbf{Approved}: true.
\item
  \textbf{Approval count}: 3.
\item
  \textbf{Required roles present}: classification\_reviewer, originator,
  release\_authority.
\item
  \textbf{Findings}: none.
\end{itemize}

The \breaktt{final\_release\_recommended} flag is true only when the
packet is releasable, the review gate is approved, and no blocking
warnings remain.
\end{frame}

\begin{frame}[allowframebreaks]{Visual Proof Matrix}
\protect\phantomsection\label{visual-proof-matrix}
The development proof matrix produces sixteen base PDFs (four redaction
styles × four backgrounds) and, when steganography is enabled, sixteen
companion steganography PDFs with hash manifests. Each variant records:

\begin{itemize}
\tightlist
\item
  Base PDF filename, byte size, and SHA-256.
\item
  Steganography PDF filename, byte size, and SHA-256.
\item
  Hash manifest filename and SHA-256.
\item
  Kmyth sidecar count and filenames (when Kmyth is available).
\end{itemize}
\end{frame}

\begin{frame}[fragile,allowframebreaks]{Kmyth TPM Sidecar Production}
\protect\phantomsection\label{kmyth-tpm-sidecar-production}
When Kmyth tools are runnable and a TPM backend is available, each
variant produces two \texttt{.ski} sidecars:

\begin{enumerate}
\tightlist
\item
  \texttt{\{variant\_id\}.hashes.json.ski} --- the hash manifest sealed
  against the TPM storage hierarchy.
\item
  \texttt{\{variant\_id\}\_steganography.pdf.ski} --- the steganography
  PDF sealed against the TPM storage hierarchy.
\end{enumerate}

The \texttt{.ski} files are ASCII-armored and contain PCR selections,
policy or-values, the storage key public area, and the sealed data
object. Unsealing requires the same TPM and platform configuration.

In the verified run, all sixteen variants produced both sidecars,
yielding thirty-two \texttt{.ski} files total. The kmyth-seal binary's
FlushContext patch ensured that consecutive seal invocations did not
exhaust the swtpm transient object slots.
\end{frame}

\begin{frame}[fragile,allowframebreaks]{Residual Risk Detection}
\protect\phantomsection\label{residual-risk-detection}
The residual-risk detector scans sanitized text for common
public-release leaks:

{\def\LTcaptype{none} % do not increment counter
\begin{longtable}[]{@{}
  >{\raggedright\arraybackslash}p{(\linewidth - 2\tabcolsep) * \real{0.5625}}
  >{\raggedright\arraybackslash}p{(\linewidth - 2\tabcolsep) * \real{0.4375}}@{}}
\toprule\noalign{}
\begin{minipage}[b]{\linewidth}\raggedright
Pattern
\end{minipage} & \begin{minipage}[b]{\linewidth}\raggedright
Regex
\end{minipage} \\
\midrule\noalign{}
\endhead
Email address &
\texttt{\textbackslash{}b{[}A-Za-z0-9.\_\%+-{]}+@{[}A-Za-z0-9.-{]}+\textbackslash{}.{[}A-Za-z{]}\{2,\}\textbackslash{}b} \\
IPv4 address &
\texttt{\textbackslash{}b(?:\textbackslash{}d\{1,3\}\textbackslash{}.)\{3\}\textbackslash{}d\{1,3\}\textbackslash{}b} \\
Coordinate pair &
\texttt{\textbackslash{}b-?\textbackslash{}d\{1,2\}\textbackslash{}.\textbackslash{}d\{3,\},\textbackslash{}s*-?\textbackslash{}d\{1,3\}\textbackslash{}.\textbackslash{}d\{3,\}\textbackslash{}b} \\
Controlled dissemination &
\texttt{\textbackslash{}b(?:NOFORN\textbackslash{}\textbar{}ORCON\textbackslash{}\textbar{}REL\textbackslash{}s+TO)\textbackslash{}b} \\
Collection discipline &
\texttt{\textbackslash{}b(?:HUMINT\textbackslash{}\textbar{}SIGINT\textbackslash{}\textbar{}IMINT\textbackslash{}\textbar{}MASINT\textbackslash{}\textbar{}OSINT)\textbackslash{}b} \\
Compartment marker &
\texttt{\textbackslash{}b(?:SCI\textbackslash{}\textbar{}TS\_SCI\textbackslash{}\textbar{}TOP\textbackslash{}s+SECRET)\textbackslash{}b} \\
Sensitive markers & \texttt{HUMINT}, \texttt{SIGINT}, \texttt{source},
\texttt{selector}, \texttt{location}, \texttt{2026-} \\
\bottomrule\noalign{}
\end{longtable}
}

Each detected marker generates a warning finding. The mosaic risk score
aggregates residual markers across all segments.
\end{frame}

\section{Discussion}\label{discussion}

\begin{frame}[allowframebreaks]{Discussion}
The exemplar demonstrates that disclosure control can be decomposed into
orthogonal concerns: text-level audit, visual presentation,
steganographic provenance, and hardware-backed sealing. Each concern is
independently configurable and verifiable.

The visual proof matrix confirms that blackout, whiteout, grayout, and
blur treatments produce equivalent source-safe outputs---only the visual
token differs. The steganography layer adds provenance without altering
the redaction decisions. Kmyth TPM sealing adds a hardware binding that
ensures sealed sidecars can only be unsealed on the same platform
configuration.

The mssim-to-swtpm protocol proxy is a necessary bridge on macOS, where
no hardware TPM exists. The proxy handles three protocol mismatches: (1)
platform commands use different command numbers, (2) the data channel
wraps TPM commands in a 9-byte mssim header, and (3) swtpm's transient
object slots are limited and must be flushed between seal invocations.

The FlushContext patch to kmyth-seal is a critical fix for batch sealing
workflows. Without it, the second kmyth-seal invocation fails with ``out
of memory for object contexts'' because the storage key from the first
invocation remains loaded in the TPM's transient object table.
\end{frame}

\begin{frame}[allowframebreaks]{Limitations}
\protect\phantomsection\label{limitations}
\begin{itemize}
\tightlist
\item
  The fixture data is invented; real release packets may require
  additional classification taxonomies and review-role policies.
\item
  The swtpm software TPM does not provide hardware-level tamper
  resistance; production deployments should use a hardware TPM.
\item
  The mssim proxy adds latency to each TPM command; batch sealing of
  thirty-two sidecars takes approximately thirty seconds.
\item
  PDF password encryption is optional and uses AES-256; the password is
  not stored in the variant matrix.
\end{itemize}
\end{frame}

\end{document}
